
\subsection{Exact comparison of all remaining path choices}

\begin{result}{B.31}{Proved}{Path choices, transported relations and an equivariant nonuniqueness certificate}
Fix the $542$ objects, the $1620$-arrow graph of B.25, its
sixty-arrow core and its $541$-edge spanning tree.  The
relative-homotopy path assignments that agree on the core
and tree form an explicitly coordinatized torsor
\[
                  \Gamma_{0,5}^{\,1068}.                       \tag{B.109}
\]
For every such assignment the formulas below give an
invertible change of generators, its exact inverse, all
$1080$ defining relators, and the comparison to the actual
fundamental groupoid.  They preserve the minimum relator
count proved in B.26.  Every assigned path has an explicit
factorization through the fixed relative-dimension-one
paths by B.21--B.25's word algorithm.

Changes in an endpoint or pin naming are covered as well:
after choosing the stated object and tree transports, the
same formulas compare any connected $542$-object,
$1620$-arrow incidence graph containing the fixed core.
The comparison is a functor on the typed arrows, not an
identification of their names or endpoints.

Moreover, the incidence sets, their $\mathfrak S_5$ action,
all displayed fractional-twist arrows, and the endpoint
plumbing germs do not by themselves select the remaining
path classes.  An actual second equivariant realization,
changing exactly the $120$ arrows $f_{\vec Q,a}$ within the
same four-point-cover families, is computed in
\texttt{b\_path\_change.py}.  It preserves that listed data,
has nonzero homology difference at every changed source,
and has all $1080$ relators transported explicitly.
The assertion here concerns this stated data.  An additional
specified marking-forgetting equation is preserved exactly
when its transported word equation holds; the criterion is
given at the end of the proof.
\end{result}
\begin{proof}
We use traversal order for path products.  Let $\mathcal F$
be the free groupoid on the full graph and
\[
 F:\mathcal F\longrightarrow\Pi_1(M_{0,5};V)
\]
the complete realization B.25.  Choose its core basepoint
$o$ and let $T_v:o\to v$ be the recorded tree path.
The union of the core and this tree has
$60+541-49=552$ edges and free rank $11$.
The other $1620-552=1068$ edges account for (B.109).
By B.21, the images of the eleven core cycles generate the
entire group $\Gamma_{0,5}$.

Fix once and for all a word function
\[
 \kappa:\Gamma_{0,5}\longrightarrow
       \{\text{based words in the sixty core arrows}\}.         \tag{B.110}
\]
This is specified by the braid normalization and combing
of B.18, the reduced $F_3\rtimes F_2$ normal form of B.19,
and the five inverse substitutions in B.21.  It satisfies
$F(\kappa(g))=g$.  A choice of reduced normal form makes it
one word function; it is not asserted to be a group
homomorphism into the free groupoid.  Its value at the
identity is the empty word.

For an alternative assignment $Q$ agreeing with $F$ on
the core and tree, define for an unfixed arrow $e:s\to t$
\[
 \delta_e=[,T_s\,Q(e)\,F(e)^{-1}\,T_s^{-1},]
           \in\Gamma_{0,5},\qquad
 L_e=T_s^{-1}\kappa(\delta_e)T_s.                              \tag{B.111}
\]
This is an actual based loop and its unique relative
homotopy class.  Conversely, for every tuple $(\delta_e)$,
concatenation with the displayed loops gives exactly one
path assignment $Q$ with those classes.  This proves the
torsor assertion, including its coordinate map.

Define the free-groupoid endofunctor
\[
 \Phi(e)=L_e e\quad\text{on the unfixed edges},\qquad
 \Phi(e)=e\quad\text{on the core and tree}.                     \tag{B.112}
\]
All the $L_e$ use only fixed edges.  Thus replacing $L_e$
by $L_e^{-1}$ in (B.112) gives its inverse, on the nose,
on every generator.  Its source and target are correct
because $L_e$ is based at $s$.  Equation (B.111) gives
\[
                         Q=F\Phi.                             \tag{B.113}
\]
Let $\mathcal R$ be the complete $1080$-relator list of
B.25.  The complete list for $Q$ is therefore
\[
                         \Phi^{-1}(\mathcal R).                \tag{B.114}
\]
Indeed $\ker Q=\Phi^{-1}(\ker F)$ as normal subgroupoids,
and $\ker F$ is normally generated by $\mathcal R$.
The induced comparison and inverse between the two
presented groupoids are exactly $\Phi,\Phi^{-1}$.
No relation has been dropped or guessed.  The fixed graph
still has rank $1079$, so B.26 gives the same lower bound
$1080$; (B.114) attains it.

The same construction supplies a comparison when a pin
choice changes an incidence target.  Here are the typed
maps, including the tree transports.  Let $\mathcal F'$
be the free groupoid on the other graph, with the same
vertex labels after a stated bijection, and choose a tree
$T'_v$ extending the same core tree.  Both graphs have
$1079$ non-tree generators.  Match their eleven fixed
core cycles, and choose a bijection $\beta$ between the
other $1068$ non-tree edges, using their recorded order.
Write
\[
 y_e=T_s eT_t^{-1},\qquad
 y'_{e'}=T'_s e'(T'_t)^{-1}
\]
for the normalized based generators.  On the alternative
graph set
\[
 \delta_{e'}=Q(y'_{e'})F(y_{\beta(e')})^{-1},\qquad
 \Phi(y'_{e'})=\kappa(\delta_{e'})y_{\beta(e')},\qquad
 \Phi(T'_v)=T_v.                                               \tag{B.115}
\]
The core cycles are fixed.  The based free-group maps in
(B.115) are invertible: invert $\beta$ and replace the
fixed core prefixes by their inverses.  The unique
transport paths in the two trees then give an inverse
on the whole free groupoids.  In particular (B.115)
means the following actual word for an arrow $e':s\to t$:
\[
 \Phi(e')=T_s^{-1}\Phi(y'_{e'})T_t.                            \tag{B.116}
\]
Define $\eta_v=F(T_v)^{-1}Q(T'_v)$.  These are explicit
paths from the reference objects to the alternative
objects, and direct substitution gives
\[
 F\Phi(e')=\eta_s\,Q(e')\,\eta_t^{-1}.                         \tag{B.117}
\]
Thus the endpoint transports give the natural isomorphism
to the actual alternative realization.  Pulling back the
relators by (B.115) again gives exactly its kernel and
$1080$ relations.  This proves the assertion for endpoint
and pin renamings as well as for path changes.

For the relative-dimension-one factorization, an arrow
$e:s\to t$ in any such realization has the based class
$[T_sQ(e)T_t^{-1}]$.  Hence its class is exactly
\[
 T_s^{-1}\kappa([T_sQ(e)T_t^{-1}])T_t.                          \tag{B.118}
\]
The core paths and the tree paths on the right already
have the explicit four-point covering or retained-circle
factorizations of B.21--B.25.  Combing computes the middle
word from any supplied path braid.  Thus the formula
factors the actual arrow, with its endpoints and its
chosen class retained.  If both original assignments
are equivariant, the comparison after imposing the
relators is equivariant: applying a label permutation
to either side of (B.113) gives the same geometric path,
and the geometric comparison is faithful on each
Hom-set.  This conclusion does not require the auxiliary
spanning tree or the word function $\kappa$ to be invariant.

Here is the promised nonuniqueness example.  At a
one-node source use the family
\[
 V(q)=(\infty,-q,q,-1,1),\qquad q\ne0,\pm1,
\]
from B.23.  Before each $f_{\vec Q,a}$ insert a full
positive meridian of $q=0$, based at its source value
$q=i\varepsilon$.  One explicit representative lets
$r=q/i$ follow the square
\[
 \varepsilon\longrightarrow i\varepsilon
   \longrightarrow-\varepsilon\longrightarrow-i\varepsilon
   \longrightarrow\varepsilon,
 \qquad\varepsilon=1/10.                                     \tag{B.119}
\]
It avoids $0,\pm1$ and is homotopic to the positive circle
around zero.  The five-point path is still the image of
the same four-point double cover.  Its class is the
meridian $u_0$ of B.23, whose image is
$\sigma_2^2=A_{23}=b$ in the corresponding normalization.
Its class in $H_1(M_{0,5};\mathbb Z)=\mathbb Z^5$ is nonzero.

The full $f_{\vec Q,a}$ family is a free
$\mathfrak S_5$ orbit of $120$ arrows.  Apply all label
permutations to this change.  Equivalently, at each of
its sixty source objects insert its relabeled meridian
in both outgoing arrows.  The source involution acts
by the constant map $z\mapsto-z$ on $V(q)$, so it fixes
this meridian path as a point of labelled moduli after
relabeling; the construction is well-defined on the
source cosets.  The resulting assignment is therefore
equivariant.  No arrow in the fractional-twist families
is changed.  All $f_{\vec Q,a}$ arrows are outside the
recorded tree, as the program checks.  Transport the meridian to an interior point along the
initial subpath, insert that conjugate loop there, and
then follow the remaining original path.  This has the
same relative class as prepending the source meridian,
while agreeing with the original path on neighborhoods
of both ends.  Its endpoint germs, framing rays, and
all endpoint incidence data are therefore unchanged.

The program computes all sixty normalized meridians,
using sixteen universal homogeneous-determinant
subinterval certificates for (B.119), and checks that
each has nonzero homology.  It computes all $120$
correction words (B.111), both inverse changes on all
$1620$ generators, and every relator in (B.114).
For example, at source vertex $122$ the meridian is
$A_{23}$, and for the first changed arrow, numbered
$241$, the old and new normalized braid words are
\[
 (2,3,2,-3,-2,-3),\qquad (2,3,2,3,-2,-3),
                                                               \tag{B.120}
\]
where $i$ means $\sigma_i$.  Their difference is the
computed source meridian, not a change of endpoint.
All $1080$ transported relators are checked by exact
inverse substitution; the generator images are also
checked by the faithful Artin action.  This proves the
claimed nonuniqueness for the stated diagram data.

Finally, the source calls the $\gamma$ arrows gommages
and subsequently defines actual marking-forgetting
functors in scans~54--60.  Those additional operations
can impose equations beyond the incidence and twist
data.  Their precise compatibility test is already
contained in the construction: for each specified typed
word equation $W=1$, the alternative assignment satisfies
it exactly when
\[
                       F\Phi(W)=1.                            \tag{B.121}
\]
Normalize and comb this explicit word using B.18--B.19;
the resulting normal form is the identity exactly in
that case.  Necessity and sufficiency are (B.113), and
(B.117) gives the transported version for changed pins.
Thus every additional supplied equation is tested by
an exact map.  The loop-insertion example asserts no
preservation of an extra equation not included in its
stated data.  Equations (B.111)--(B.118) supply the full
choice-independent comparison for every actual
realization, including those satisfying the extra
marking functor equations.
\end{proof}
